\documentclass[11pt]{article}
\usepackage{amsmath}
\usepackage{booktabs}
\usepackage{listings}

\title{Comparing Two Document Conversion Workflows}
% TODO: Add author name and institutional affiliation.
\author{}
% TODO: Add the report date if required.
\date{}

\begin{document}
\maketitle

\noindent\textbf{Status:} Fictional demonstration report; not real research findings.

\section{Objective}
Compare the time required to prepare editable LaTeX drafts from Markdown
notes using Workflow A and Workflow B.

\section{Methods}
Each workflow was applied to the same 6 fictional Markdown reports. Each
report contained one table and one equation. Each workflow was run once
per report. Preparation time included manual correction of the generated
LaTeX. Hardware and software versions were not recorded.
% TODO: Add hardware and software details if available.

\section{Results}
Table~\ref{tab:workflow-results} presents the values supplied by the report
author.

\begin{table}[htbp]
\centering
\caption{Reported results for the fictional document conversion demonstration.}
\label{tab:workflow-results}
\begin{tabular}{lrrr}
\toprule
Workflow & Reports & Mean preparation & Reports requiring \\
         & processed & time (min) & manual correction \\
\midrule
A & 6 & 18.5 & 4 \\
B & 6 & 12.0 & 2 \\
\bottomrule
\end{tabular}
\end{table}

\subsection{Reported Equation}
Mean preparation time was defined as
\begin{equation}
\bar{t} = \frac{1}{n}\sum_{i=1}^{n} t_i
\label{eq:mean-time}
\end{equation}
where \(t_i\) is preparation time for report \(i\), measured in minutes,
and \(n\) is the number of reports.

\subsection{Code Excerpt}
This Python excerpt labels the two workflows; it does not perform analysis.
\begin{lstlisting}[language=Python,breaklines=true,basicstyle=\ttfamily\small]
workflow_labels = ["A", "B"]
print("Workflow labels:", workflow_labels)
\end{lstlisting}

\section{Discussion}
Workflow B had a lower reported mean preparation time in this demonstration.
No significance tests, confidence intervals, or effect sizes were computed.
These observations do not establish a general performance advantage.

% TODO: Supply a source for the following background statement.
Structured authoring workflows can support reproducible document preparation.
\textbf{[Citation needed; no source supplied.]}

\section{Limitations}
Only 6 fictional reports were included. Each workflow was run once per
report. Hardware and software versions were not recorded. Manual
corrections were included in preparation time but were not categorized
by error type.

\section{Missing Materials}
Author name and affiliation were not supplied. A workflow diagram was
mentioned in the planning notes, but no figure file was supplied. No
bibliography was supplied.
% TODO: Add the workflow diagram only when a figure file is supplied.
% TODO: Add references only when bibliographic information is supplied.

\end{document}
